% !TeX root = ../main.tex
% !TeX encoding = utf8
\section{NGHIÊN CỨU MÔ PHỎNG MONTE CARLO BẰNG R}

\subsection{Thiết kế mô phỏng}

Mô phỏng đánh giá coverage và độ dài trung bình của Normal, Basic và
Percentile CI với mức tin cậy danh nghĩa $95\%$. Mỗi trong bốn kịch bản
được lặp $K=10{,}000$ lần; mỗi lần lặp tạo $B=1{,}999$ mẫu Bootstrap.
Các kịch bản gồm tổng thể chuẩn $N(0,1)$ và mũ $\operatorname{Exp}(1)$,
đều ở hai cỡ mẫu $n=12$ và $n=40$. Seed được cố định ở $20260925$ để có
thể tái lập kết quả.

Coverage là tỷ lệ các khoảng chứa tham số thật; độ dài trung bình là trung
bình của $U_k-L_k$ qua các lần lặp. Sai số chuẩn Monte Carlo của coverage
được ước lượng bởi
\[
\widehat{MCSE}=\sqrt{\widehat{CP}(1-\widehat{CP})/K}.
\]
Trong ví dụ liệt kê chính xác ở Phần~\ref{sec:examples}, phân vị dùng
quy ước type~1; mô phỏng dùng phân vị nội suy type~7 của R. Tệp
\texttt{run-info.txt} trong thư mục kết quả ghi lại các thiết lập này.

Đoạn mã sau thể hiện phần tính các cận từ một mẫu Bootstrap. Vòng lặp đầy
đủ qua bốn kịch bản, tổng hợp chỉ số và xuất kết quả được triển khai trong
mã mô phỏng R.

\begin{lstlisting}[style=Rstyle, caption={Tính các khoảng tin cậy}]
means <- colMeans(matrix(
  x[sample.int(n, n * B, replace = TRUE)], nrow = n))
se <- sd(means)
q <- quantile(means, c(alpha / 2, 1 - alpha / 2), type = 7)
ci <- rbind(
  Normal = c(theta_hat - z * se, theta_hat + z * se),
  Basic = c(2 * theta_hat - q[2], 2 * theta_hat - q[1]),
  Percentile = q
)
\end{lstlisting}

\subsection{Kết quả mô phỏng}

Bảng~\ref{tab:simulation-results} tổng hợp kết quả từ lần chạy có seed cố
định. Coverage và MCSE được trình bày theo phần trăm; AL giữ nguyên đơn vị
của thống kê trung bình.

\begin{table}[htbp]
\centering
\scriptsize
\renewcommand{\arraystretch}{1.12}
\setlength{\tabcolsep}{4pt}
\caption{Coverage, độ dài khoảng trung bình và MCSE trong mô phỏng ($K=10{,}000$, $B=1{,}999$)}
\label{tab:simulation-results}
\begin{tabular}{llrrr}
\toprule
Phân phối & Phương pháp & Coverage (\%) & AL & MCSE (điểm \%) \\
\midrule
\multirow{3}{*}{$N(0,1),\ n=12$}
& Normal & 91.05 & 1.0585 & 0.29 \\
& Basic & 91.01 & 1.0528 & 0.29 \\
& Percentile & 90.74 & 1.0528 & 0.29 \\
\addlinespace
\multirow{3}{*}{$N(0,1),\ n=40$}
& Normal & 93.91 & 0.6081 & 0.24 \\
& Basic & 93.89 & 0.6065 & 0.24 \\
& Percentile & 93.79 & 0.6065 & 0.24 \\
\addlinespace
\multirow{3}{*}{$\operatorname{Exp}(1),\ n=12$}
& Normal & 87.06 & 1.0153 & 0.34 \\
& Basic & 85.47 & 1.0002 & 0.35 \\
& Percentile & 87.59 & 1.0002 & 0.33 \\
\addlinespace
\multirow{3}{*}{$\operatorname{Exp}(1),\ n=40$}
& Normal & 91.92 & 0.5971 & 0.27 \\
& Basic & 90.99 & 0.5939 & 0.29 \\
& Percentile & 92.21 & 0.5939 & 0.27 \\
\bottomrule
\end{tabular}
\end{table}

\begin{figure}[htbp]
\centering
\includegraphics[width=0.95\textwidth]{simulation-coverage.pdf}
\caption{Coverage của ba phương pháp trong bốn kịch bản; đường đứt biểu thị mức danh nghĩa $95\%$.}
\label{fig:simulation-coverage}
\end{figure}

\begin{figure}[htbp]
\centering
\includegraphics[width=0.95\textwidth]{simulation-average-length.pdf}
\caption{Độ dài khoảng trung bình của ba phương pháp trong bốn kịch bản.}
\label{fig:simulation-average-length}
\end{figure}

Trong tổng thể chuẩn, coverage của ba phương pháp khá gần nhau và tăng từ
khoảng $91\%$ ở $n=12$ lên khoảng $94\%$ ở $n=40$. Với tổng thể mũ, ở cả
hai cỡ mẫu, Basic có coverage thấp hơn Normal trong lần mô phỏng này;
Percentile có coverage cao nhất trong ba phương pháp ở cả $n=12$ và $n=40$.
AL của Basic và Percentile trùng nhau trong từng kịch bản do hai khoảng
được phản chiếu từ cùng cặp phân vị. Các kết quả này mô tả đúng bốn kịch
bản đã xét, không khẳng định thứ hạng phương pháp cho mọi phân phối.

